\section{Signed defects and two routes to reassembly}
\label{sec:tpc60-signed-defect}

The Gram formulation is optimal for a coefficient-uniform theorem.  For a
distinguished physical coefficient vector, the exact remaining task is
often clearer as a signed cross-term estimate.

\subsection{The pointwise signed defect}

Define
\begin{equation}
 \mathfrak X(z)
 :=2\operatorname{Re}\langle Az,Mz\rangle+\|Mz\|^2.
 \label{eq:tpc60-signed-defect}
\end{equation}
Then
\begin{equation}
 \boxed{
 \|Bz\|^2=\|Az\|^2+\mathfrak X(z).}
 \label{eq:tpc60-signed-defect-identity}
\end{equation}
For \(Az\ne0\), this becomes
\begin{equation}
 \boxed{
 \frac{\|Bz\|^2}{\|Az\|^2}
 =1+\frac{\mathfrak X(z)}{\|Az\|^2}.}
 \label{eq:tpc60-signed-ratio}
\end{equation}
Consequently a desired pointwise transfer
\begin{equation}
 \|Bz\|^2\ge\delta_*\|Az\|^2
 \label{eq:tpc60-desired-pointwise-transfer}
\end{equation}
is exactly equivalent to
\begin{equation}
 \boxed{
 \mathfrak X(z)
 \ge-(1-\delta_*)\|Az\|^2.}
 \label{eq:tpc60-exact-signed-target}
\end{equation}
This is neither a positivity statement about \(\|Mz\|^2\) nor an unsigned
estimate for the atomic missing mass.  The cross term is part of the
physical problem.

In the prime-resolved notation of
\cref{prop:tpc60-same-prime-separation}, \(\mathfrak X\) is exactly the
right side of \eqref{eq:tpc60-exact-prime-alias-defect}.  Hence its negative
part can arise only from base--missing and unequal-prime native aliasing.
The same-prime orthogonality is a genuine simplification, but it leaves the
two surviving sums uncontrolled.

\subsection{The nuisance-optimized defect}

For \(s\in S\), define the best defect obtainable by moving inside its
\(A\)-fiber:
\begin{equation}
 \mathfrak X_{\rm sh}(s)
 :=\inf_{k\in N}
 \left(
  \|B(s+k)\|^2-\|A_Ss\|^2
 \right).
 \label{eq:tpc60-shorted-signed-defect}
\end{equation}

\begin{proposition}[Shorted defect operator]
\label{prop:tpc60-shorted-defect-operator}
With \(Q\) and \(R_S\) as in
\cref{sec:tpc60-shorted-gram},
\begin{equation}
 \boxed{
 \mathfrak X_{\rm sh}(s)
 =\langle(Q-R_S)s,s\rangle.}
 \label{eq:tpc60-shorted-defect-form}
\end{equation}
If
\begin{equation}
 \Theta_{\rm reasm}
 :=R_S^{-1/2}(Q-R_S)R_S^{-1/2},
 \label{eq:tpc60-Theta-reassembly}
\end{equation}
then
\begin{equation}
 \boxed{
 \delta(A,M)
 =1+\lambda_{\min}(\Theta_{\rm reasm}).}
 \label{eq:tpc60-delta-Theta}
\end{equation}
\end{proposition}

\begin{proof}
The least-squares identity
\eqref{eq:tpc60-distance-identity} gives
\[
 \inf_{k\in N}\|B(s+k)\|^2
 =\langle Qs,s\rangle.
\]
Subtracting \(\|A_Ss\|^2=\langle R_Ss,s\rangle\) proves
\eqref{eq:tpc60-shorted-defect-form}.  Conjugating by \(R_S^{-1/2}\) and
using \cref{thm:tpc60-shorted-generalized-eigenvalue} proves
\eqref{eq:tpc60-delta-Theta}.
\end{proof}

Thus a coefficient-uniform signed route must control the negative spectrum
of the \emph{shorted} defect.  Bounding the unshorted cross term only on the
canonical sector \(S\) is insufficient when \(B(N)\ne0\).

\subsection{An angle certificate}

For \(Az\ne0\), put
\begin{equation}
 t(z):=\frac{\|Mz\|}{\|Az\|},
 \label{eq:tpc60-size-ratio-t}
\end{equation}
and, when \(Mz\ne0\), define the destructive cosine
\begin{equation}
 \rho(z)
 :=-\frac{\operatorname{Re}\langle Az,Mz\rangle}
          {\|Az\|\,\|Mz\|}.
 \label{eq:tpc60-destructive-cosine}
\end{equation}
Set \(\rho(z)=0\) if \(Mz=0\).  Then \(-1\le\rho(z)\le1\), and
\begin{equation}
 \frac{\|Bz\|^2}{\|Az\|^2}
 =1+t(z)^2-2\rho(z)t(z).
 \label{eq:tpc60-angle-size-identity}
\end{equation}

\begin{proposition}[Uniform angle certificate]
\label{prop:tpc60-uniform-angle-certificate}
Suppose there is \(0\le\rho_0<1\) such that
\begin{equation}
 \rho(z)\le\rho_0
 \qquad\text{whenever }Az\ne0.
 \label{eq:tpc60-uniform-angle-hypothesis}
\end{equation}
Then
\begin{equation}
 \boxed{
 \delta(A,M)\ge1-\rho_0^2>0.}
 \label{eq:tpc60-uniform-angle-bound}
\end{equation}
In particular, if
\(\operatorname{Re}\langle Az,Mz\rangle\ge0\) for all \(z\), then
\(\delta(A,M)\ge1\).
\end{proposition}

\begin{proof}
If \(\rho(z)\ge0\), complete the square:
\[
 1+t^2-2\rho t=(t-\rho)^2+1-\rho^2
 \ge1-\rho_0^2.
\]
If \(\rho(z)<0\), the left side is at least one.  Take the infimum over
\(z\).
\end{proof}

This angle route may succeed even when the missing norm is not small.  It
is also substantially more arithmetic: one must retain the phases in the
base--missing and unequal-prime alias sums.  The mass--degree route in the
next section is cruder, but it replaces phase control by a quantitative
smallness condition.
